\subsection{One historical scalar gives no direction}

Even if the denominator problem is repaired, one chosen policy reveals no
counterfactual cost for another policy.

\begin{proposition}[One-signal nonidentifiability]
\label{prop:nonidentify}
Fix two candidate future policies.  Over unrestricted repeated-task cost
sequences, no deterministic rule that sees the same historical revisit
record in two environments can choose a bounded-competitive future policy in
both.  For a randomized rule and every $M>1$, one of two indistinguishable
environments has expected competitive ratio at least $(M+1)/2$.
\end{proposition}

\begin{proof}
Let the complete history, including $\varrho_{\rm rev}$, be identical.  In
environment one, assign future policy costs $(1,M)$; in environment two,
assign $(M,1)$.  A deterministic rule makes the same choice in both and pays
$M$ in one of them while the best policy costs one.  If it chooses the first
policy with probability $p$, its expected costs are
$p+(1-p)M$ and $pM+(1-p)$.  Their maximum is at least their average,
$(M+1)/2$.
\end{proof}

This is a decision-theoretic result, not an SOR graph construction.  Its role
is precise: a guarantee from the historical scalar requires an additional
assumption such as stationary tasks, a relation between contexts and costs,
counterfactual probes, or a safe portfolio that runs more than one policy.
Keying memory only by graph size $n$ supplies none of these assumptions;
nonisomorphic graphs, different seeds, and different gates can share $n$.
